\documentclass[11pt]{article}
\usepackage[margin=1in]{geometry}
\usepackage{amsmath,amssymb,amsthm,mathtools}
\usepackage[T1]{fontenc}
\usepackage{hyperref}
\newtheorem{theorem}{Theorem}
\newtheorem{lemma}[theorem]{Lemma}
\newtheorem{proposition}[theorem]{Proposition}
\newtheorem{corollary}[theorem]{Corollary}
\newcommand{\R}{\mathbb R}
\newcommand{\T}{\mathcal T}
\newcommand{\Crit}{\operatorname{Crit}}
\newcommand{\supp}{\operatorname{supp}}
\newcommand{\ip}[2]{\langle #1,#2\rangle}
\newcommand{\norm}[1]{\lVert #1\rVert}
\title{Hamiltonian transport with a varying tube class at infinity\\
\large Closing the implication in manuscript 340}
\author{Campaign 9000340: Engineer proof for independent review}
\date{5 November 2026}
\begin{document}
\maketitle

\begin{abstract}
We prove the missing transport implication by the direct small-step
construction. All four differentiated cutoff equations are retained.
A block right inverse proves full Morse-family regularity without
inverting a fiber Hessian. Explicit uniform value and full-vertical
derivative bounds permit application of the already reviewed smoothing
lemma. The time subdivision depends only on the Hamiltonian and a
spatial buffer, not on the families subsequently constructed. Consequently
the fixed \(S^9\) slice retains its actual, possibly nonconstant, stable
tube class. The accepted realization and one-critical obstruction then
give the claimed non-Hamiltonian-isotopy conclusion.
This is a prose verification using accepted campaign theorems, supported
by exact algebraic replay. It is not a formal verification of topology
or a certification of novelty.
\end{abstract}

\section{Precise inputs and conclusion}

Our conventions are
\[
 \lambda=p\,dq,\qquad \omega=d\lambda,\qquad
 \iota_{X_H}\omega=-dH.
 \tag{1}
\]
Thus \(\dot q=H_p\) and \(\dot p=-H_q\).
A Morse family \(F(x,w)\) means that its \emph{full} equation
\(F_w:U\times\R^r\to\R^r\) is transverse to zero.
This does not require the Hessian \(F_{ww}\) to be invertible.
Its generating map is \(\Gamma_F(x,w)=(x,F_x(x,w))\) on
\(\Crit(F)=F_w^{-1}(0)\).

Write \(q_{k,l}(u,v)=-|u|^2+|v|^2\).
The tube model is the \(C^1\) component of its angular restriction
among functions with zero a regular value, stabilized by \(q_{1,1}\).
Homogeneous data extend by \(g(ru)=r^2h(u)\), including their
\(C^1\) extension at zero. The spherical class is that of the
negative domains, exactly as in the two accepted proofs.

We use the following existing mathematical inputs, without rerunning
their unchanged calculation certificates.
\begin{enumerate}
\item \textbf{Accepted realization (TR-existence).}
Theorem 1 of the accepted realization proof%
\footnote{\path{research/tube_realization_audit/proof.tex}}
supplies some \(N=2m\), \(m\geq20\), and a smooth Morse family
\(F:Q\times\R^{2m}\to\R\), \(Q=S^9\times S^{N-1}\).
Its compact critical locus is diffeomorphic to \(Q\), and
\(\Gamma_F\) is a diffeomorphism onto a closed exact embedded
Lagrangian \(L\). Outside a common ball, \(F(b,v,w)=g_b(w)\).
On each fixed \(S^9\times\{v_0\}\), these data represent the same
nonzero stable tube class \(a\), with spherical class zero.
The independent review is
\path{handoffs/eabe9d4c259e/round-0001.json} of the assigned project.
\item \textbf{Accepted obstruction (MS-one-critical).}
Theorem 1 of \path{research/one_critical_audit/proof.tex}
says that a smooth family on \(S^9\times\R^{2p}\), \(p\geq20\),
with uniform smooth homogeneous \(T_{p,p}\) infinity data, one
smooth nondegenerate critical section, and trivial stable spherical
class has trivial stable tube class.
Its independent review is
\path{handoffs/c4428bac0f24/round-0001.json}.
\item \textbf{Accepted quantitative smoothing.}
The lemma labelled \texttt{lem:smoothing} in that same proof applies
to \(G_0(x,w,y)=g_x(w)+q'(y)\), with \(q'\) fixed nondegenerate,
\(|d_wg_x|\geq c|w|\), and a smooth \(\widehat F\) satisfying
\[
 |\widehat F-G_0|\leq M_0,\qquad
 |d_{w,y}(\widehat F-G_0)|\leq M_1
 \tag{2}
 \]
uniformly. It changes \(\widehat F\) only outside an arbitrarily
large common ball, introduces no critical points, and produces
exact smooth homogeneous data at infinity regularly \(C^1\)-homotopic
to \(G_0\). Its proof explicitly permits any parameter set on which
the bounds are uniform; it is not restricted to sphere parameters.
\end{enumerate}
These are internal, independently reviewed inputs already present in
the ledger. Their external topological assumptions remain those of
the accepted proofs. No new literature theorem is assumed here.
In particular, we do not invoke AIS Lemma 3.23, Sikorav, or a general
isotopy-invariance theorem. There is no unexamined Legendrian-lift
convention hidden in the argument.

\begin{theorem}\label{thm:transport}
Suppose the \(L\) and \(F\) of the accepted realization admit a
compactly supported Hamiltonian isotopy taking \(L\) to \(Q_0\).
For any fixed \(v_0\in S^{N-1}\) there then exist a finite
\(p\geq20\) and a smooth family \(f:S^9\times\R^{2p}\to\R\)
with one smoothly varying nondegenerate critical point in each
fiber, uniform smooth homogeneous \(T_{p,p}\) infinity data,
and the same free stable tube and spherical classes as \(g_b\).
In particular its tube class is nonzero and its spherical class is zero.
\end{theorem}

\begin{corollary}\label{cor:main}
Composing Theorem~\ref{thm:transport} with the accepted obstruction,
the closed exact embedded \(L\) of TR-existence is not compactly
supported Hamiltonian isotopic to the zero section.
Thus, subject to the explicitly accepted construction and obstruction,
the main mathematical conclusion of manuscript 340 follows.
\end{corollary}

\section{A reverse isotopy on a Euclidean cylinder}

Put \(d=10+N\). For \(0<\nu<1\), identify a bounded neighborhood
\(U\subset\R^{10}\times\R^N\) with \(Q\times(-\nu,\nu)^2\) by
\[
 ((b,v),\alpha)\longmapsto ((1+\alpha_1)b,(1+\alpha_2)v).
\]
This chart, its radial coordinates, and its projection \(r_Q\)
are smooth on a neighborhood of the compact closure.
The cotangent coordinates satisfy
\(\lambda_U=\lambda_Q+p_\alpha\,d\alpha\).
The family \(F^{(0)}(x,w)=F(r_Q(x),w)\) generates precisely
\[
 \mathcal L_0=\{(q,\alpha;p_q,0):(q,p_q)\in L\}.
 \tag{3}
\]
It is a Morse family, and its generating parametrization is a
diffeomorphism. Indeed the old full-gradient surjectivity persists
under pullback by the submersion \(r_Q\); the additional critical-locus
coordinates are exactly \(\alpha\).
All estimates are uniform up to \(\overline U\).

For completeness, if the hypothesized flow \(\varphi_t\), generated
by \(H_t\), is given from \(Q_0\) to \(L\), the reverse path is
\[
 \psi_t=\varphi_{1-t}\varphi_1^{-1},\qquad H_t^Q=-H_{1-t}.
 \tag{4}
\]
Differentiation gives
\(\dot\psi_t=-X_{H_{1-t}}\circ\psi_t\).
There is no conjugating factor: \(\varphi_1^{-1}\) is fixed on the
\emph{right}. This differs from differentiating \(\varphi_t^{-1}\).

Choose a smooth \(0\leq\zeta\leq1\), compactly supported in the
radial square and equal to one on a smaller square.
Set
\[
 M=\sup|H_t^Q|\,\sup|d\zeta|.
\]
Take a smooth compactly supported \(0\leq\xi\leq1\) on the radial
momentum space, equal to one for \(|p_\alpha|\leq M+1\), and extend
\[
 \widetilde H_t=H_t^Q(q,p_q)\zeta(\alpha)\xi(p_\alpha)
 \tag{5}
\]
by zero from \(T^*U\) to \(T^*\R^d\).
It has one compact support for all times. Its radial equations are
\[
 \dot\alpha=H_t^Q\zeta\,d\xi,\qquad
 \dot p_\alpha=-H_t^Q\xi\,d\zeta.
 \tag{6}
\]
For every initial point of (3), \(|p_\alpha(t)|\leq Mt\);
this estimate uses only \(0\leq\xi\leq1\), so is not a circular
plateau argument. Thus \(d\xi=0\) all along the trajectory and
\(\alpha\) stays fixed. On the central radial square,
\(d\zeta=0\), \(p_\alpha=0\), and the \(Q\) motion is the reverse
isotopy. No other radial slice can enter it. Consequently, for
a smaller cylinder \(U_c\),
\[
 \widetilde\varphi_1(\mathcal L_0)\cap T^*U_c=(U_c)_0.
 \tag{7}
\]

\section{Fixing the subdivision before constructing any new family}

Choose a compact \(K\Subset T^*U\) containing all supports of
\(\widetilde H_t\) in its interior. Every partial time flow fixes
the complement of \(K\); uniqueness, also backwards in time,
prevents a point of \(K\) from flowing to a stationary point outside
\(K\). All partial flows therefore preserve \(K\) and \(T^*U\).

For a sufficiently short interval let
\(\psi(x,p)=(X,P)\) be its time increment.
Require \(\|D_xX-I\|<1/2\). For fixed \(p\), integration along a
line segment gives
\[
 |X(x_1,p)-X(x_2,p)|\geq\tfrac12|x_1-x_2|.
 \tag{8}
\]
The map \(x\mapsto X(x,p)\) is a proper local diffeomorphism,
equal to the identity outside a compact set. Its image is open and
closed in \(\R^d\); it is therefore a global diffeomorphism.
The resulting \(x=x(X,p)\) is smooth jointly in the two variables.

Cartan's formula with (1) gives
\(\mathcal L_{X_{\widetilde H}}\lambda
=d(\lambda(X_{\widetilde H})-\widetilde H)\).
For the partial flow \(\varphi_{s,a}\) from \(a\) to \(s\), put
\[
 A(x,p)=\int_a^b
       (\lambda(X_{\widetilde H_s})-\widetilde H_s)
                    (\varphi_{s,a}(x,p))\,ds.
 \tag{9}
\]
Then \(dA=P\,dX-p\,dx\) on the graph. Hence
\[
 \sigma(X,p)=A(x(X,p),p)+(x(X,p)-X)\cdot p
 \tag{10}
\]
satisfies
\[
 x=X+\sigma_p,\qquad P=p+\sigma_X,\qquad
 J:=I+\sigma_{pX}=D_Xx\ \hbox{invertible}.
 \tag{11}
\]
The function \(A\) vanishes outside \(K\), so \(\sigma\) is
compactly supported in the mixed-coordinate image of \(K\).
Choose a compact \(D\Subset U\) containing the spatial projection
of \(K\) in its interior. This one \(D\) works for \emph{every}
short interval, not just for a previously selected subdivision.
On \(U\setminus D\), \(\sigma\) and every derivative vanish for
all \(p\).

Choose \(\zeta_1\in C^\infty_c(U,[0,1])\), equal to one on an
open \(V\supset D\), and fix \(\delta>0\) so that
\[
 2\delta<\operatorname{dist}(\supp\zeta_1,\partial U).
 \tag{12}
\]
These choices depend only on the Hamiltonian support and \(U\).
Let \(L_1=\sup\|DX_{\widetilde H_t}\|\) and
\(V_1=\sup|X_{\widetilde H_t}|\), both finite.
Every increment of length at most
\[
 h=\min\{1,\ \log(3/2)/(1+L_1),\ \delta/(1+V_1)\}>0
 \tag{13}
\]
has \(\|D\psi-I\|<1/2\) and \(\|\sigma_p\|_\infty<\delta\).
The first assertion follows from
\(\|D\psi-I\|\leq e^{L_1h}-1\), and the second from displacement
bounded by \(V_1h\). Fix \(s<\infty\) with \(1/s\leq h\).
No later choice of generating-family constants changes this subdivision.

\section{One step, including every spatial-boundary term}

Assume the old smooth Morse family \(F_o:U\times\R^r\to\R\)
has a diffeomorphic generating parametrization, extends smoothly
past \(\overline U\), and equals \(g_o(X,w)\) for \(|w|\geq R_o>0\).
Assume, uniformly up to the boundary,
\[
 |(g_o)_w(X,w)|\geq c|w|,\qquad
 |(g_o)_X(X,w)|\leq C|w|^2,\qquad c>0,\quad C\geq0.
 \tag{14}
\]
The initial family has these properties by compactness, regular
homogeneous angular data, and the accepted realization.

Let \(\beta(z)\) be one on \(|z|\leq\delta\) and zero on
\(|z|\geq2\delta\). Choose \(R_1>R_o\), and a smooth radial
\(\chi(w)\), one on a neighborhood of \(|w|\leq R_1\),
zero outside a finite ball, with
\[
 |d\chi(w)|\leq\delta_1/|w|,\qquad
 \delta_1=\frac{c}{4\delta(1+C)}.
 \tag{15}
\]
Such a cutoff exists: apply a fixed smooth plateau function
\(\theta\) to \(\log(|w|/R_1)/T\). Choose
\[
 T\geq\max(1,\|\theta'\|_\infty/\delta_1),
\]
and extend by one on the inner ball.
Its support radius can be enormous, but is finite.

For \(z,p\in\R^d\), set
\[
 k=\zeta_1(X)\beta(z)\chi(w),\qquad x=X+kz,\qquad
 \widehat F(X,w,z,p)=F_o(x,w)+\sigma(X,p)-z\cdot p.
 \tag{16}
\]
When the shift is nonzero, \(X\in\supp\zeta_1\) and its length
is at most \(2\delta\). The entire segment from \(X\) to \(x\)
is inside \(U\), by (12). When the shift is zero there is no
boundary issue. This also proves smooth extension up to
\(\overline U\); all base cutoffs vanish near its boundary.

Put \(a=(F_o)_x(x,w)\). The full derivatives are
\begin{align}
 \widehat F_p&=\sigma_p-z, \tag{17}\\
 \widehat F_z&=ka+(a\cdot z)k_z-p, \tag{18}\\
 \widehat F_w&=(F_o)_w(x,w)+(a\cdot z)k_w, \tag{19}\\
 \widehat F_X&=a+(a\cdot z)k_X+\sigma_X. \tag{20}
\end{align}
In particular
\(k_z=\zeta_1\chi\,d\beta\),
\(k_w=\zeta_1\beta\,d\chi\), and
\(k_X=\beta\chi\,d\zeta_1\).
These identities hold before imposing any critical equation.

At a vertical critical point, (17) implies
\(z=\sigma_p(X,p)\), so \(|z|<\delta\) and \(\beta=1\)
on a neighborhood. If \(|w|\geq R_o\), (14)--(15) give
\[
 |\widehat F_w|
 \geq(c-C\delta\delta_1)|w|
 =c\left(\frac34+\frac1{4(1+C)}\right)|w|>0.
 \tag{21}
\]
Thus every critical point has \(|w|<R_o\), and \(\chi=1\)
on a neighborhood as well. There are exactly two local cases,
covering \(U\):
\begin{itemize}
\item On \(V\), all cutoffs equal one near a solution. The
equations are
\((F_o)_w=0\), \(p=(F_o)_x\), \(x=X+z=X+\sigma_p\).
Equation (20) gives the correct transported momentum
\(P=p+\sigma_X\).
\item On \(U\setminus D\), \(\sigma=0\) for all momenta, and the
equations are
\[
 z=0,\quad x=X,\quad (F_o)_w(X,w)=0,\quad
 p=\zeta_1(X)(F_o)_x(X,w).
 \tag{22}
\]
The generated covector is nevertheless \((F_o)_x(X,w)\), not
the possibly rescaled auxiliary variable \(p\). The spatial
cutoff term in (20) vanishes because \(z=0\).
This is exactly the identity symplectomorphism on this region.
\end{itemize}
On the overlap \(\zeta_1=1\), so the descriptions agree.

\subsection{Full regularity without a fiberwise Morse assumption}

In the first case use the total coordinates \((x,w,p,X)\).
The equations, up to their ordering, are
\[
 E_1=(F_o)_w,\qquad E_2=(F_o)_x-p,\qquad
 E_3=X+\sigma_p(X,p)-x.
 \tag{23}
\]
At a solution their differential is the block matrix
\[
 D E=
 \begin{pmatrix}
 A&C&0&0\\ B&A^{\mathsf T}&-I&0\\-I&0&S&J
 \end{pmatrix},
 \quad
 A=(F_o)_{wx},\ C=(F_o)_{ww},\
 B=(F_o)_{xx},\ S=\sigma_{pp}.
 \tag{24}
\]
Old full regularity provides \(U_0,V_0\) with
\(AU_0+CV_0=I_r\). Set \(K_0=BU_0+A^{\mathsf T}V_0\).
An explicit right inverse for (24) is
\[
 \begin{pmatrix}
 U_0&0&0\\ V_0&0&0\\ K_0&-I&0\\
 J^{-1}(U_0-SK_0)&J^{-1}S&J^{-1}
 \end{pmatrix}.
 \tag{25}
\]
Multiplication gives diagonal blocks \(I_r,I_d,I_d\) and zero
off-diagonal blocks. Neither the construction nor the identity
uses \(C^{-1}\).

In the second case, the equation \(-z=0\) first eliminates \(z\);
then (18), with coefficient \(-I\) in \(p\), eliminates \(p\).
The remaining equation is precisely the old full-gradient equation
on \((X,w)\), so is regular. This argument includes the entire
transition of \(\zeta_1\), with no assumption that its derivative
vanishes there.

\subsection{The global critical-locus correspondence}

For an old critical point \((x,w)\), let
\(p_o=(F_o)_x(x,w)\) and \((X,P)=\psi(x,p_o)\).
The correspondence and its inverse are
\[
 (x,w)\longmapsto(X,w,x-X,\zeta_1(X)p_o),\qquad
 (X,w,z,p)\longmapsto(X+kz,w).
 \tag{26}
\]
The two cases above verify that these smooth maps are inverse and
that the generating maps intertwine with \(\psi\).
Since \(\psi\) preserves \(T^*U\), there is no missing branch
entering or leaving through a spatial boundary.
All critical points obey the uniform bounds
\[
 |w|<R_o,\quad |z|<\delta,\quad
 |p|\leq\sup_{\overline U,\ |w|\leq R_o}|(F_o)_x|.
 \tag{27}
\]
Thus \(\widehat F\) is a Morse family with a diffeomorphic
parametrization of precisely \(\psi(\mathcal L_o)\).

\section{Full-vertical bounded error and the actual infinity class}

Set \(W=(w,z,p)\) and
\[
 G_0(X,W)=g_o(X,w)-z\cdot p.
 \tag{28}
\]
This is \(C^1\), homogeneous of degree two, and
\[
 |(G_0)_W|^2=|(g_o)_w|^2+|z|^2+|p|^2
 \geq\min(c^2,1)|W|^2.
 \tag{29}
\]
It need not be smooth on the \(w=0\) axis; this is exactly the
situation covered by the accepted smoothing lemma.

We check (2) for the \emph{full} vertical differential, not merely
the old \(w\) variables. The cancellation is exact:
\[
 E:=\widehat F-G_0
   =\underbrace{F_o(X+kz,w)-F_o(X,w)}_{\Delta}
       +\underbrace{F_o(X,w)-g_o(X,w)}_{R}
       +\sigma(X,p).
 \tag{30}
\]
Let \(\supp\chi\subset\{|w|\leq R_\chi\}\).
On \(\overline U\times\{|w|\leq R_\chi\}\), choose bounds
\(B_x\geq|(F_o)_x|\) and \(B_w\geq|(F_o)_w|\).
Let \(K_\chi=\|d\chi\|_\infty\), \(K_\beta=\|d\beta\|_\infty\);
let \(|R|\leq C_0\), \(|R_w|\leq C_1\); and let
\(|\sigma|\leq S_0\), \(|\sigma_p|\leq S_1\).
All are finite. The bounds for \(R\) hold because it vanishes
outside \(|w|\leq R_o\) and is \(C^1\), including at zero.

The function \(\Delta\) and each of its vertical derivatives vanish
unless both \(w\) and \(z\) are in their fixed support balls.
Using the segment inside \(U\) and differentiating (30) gives
\begin{align}
 |\Delta|&\leq2\delta B_x, \tag{31}\\
 |\Delta_w|&\leq2B_w+2\delta B_xK_\chi, \tag{32}\\
 |\Delta_z|&\leq B_x(1+2\delta K_\beta),\qquad
 \Delta_p=0. \tag{33}
\end{align}
The second terms in (32) and (33) are the differentiated cutoffs;
they cannot be omitted. By the triangle inequality for the
vertical blocks, valid explicit constants in (2) are
\begin{align}
 M_0&=2\delta B_x+C_0+S_0, \tag{34}\\
 M_1&=2B_w+2\delta B_xK_\chi+
         B_x+2\delta B_xK_\beta+C_1+S_1. \tag{35}
\end{align}
These constants are independent of \(p,z,w,X\).
We do not claim that \(E\) has compact support in \emph{all}
vertical variables: \(\Delta\) is independent of \(p\).
Bounded value and bounded full vertical derivative are what is needed.

Apply the accepted smoothing lemma, uniformly over \(\overline U\),
with \(q'(z,p)=-z\cdot p\) and with its untouched ball larger
than (27). Its endpoint \(F_n\) is smooth, exactly smooth
homogeneous at infinity, and identical to \(\widehat F\) on a
neighborhood of the entire critical locus. Thus all critical points,
the full-gradient regularity, and the generated base covectors are
unchanged. This is not just preservation of the number of zeros.
The lemma supplies the regular homogeneous homotopy from (28)
to the new infinity data. Its four-error cutoff estimate has already
been reviewed; the new assertion here is precisely that (34)--(35)
meet its hypotheses.

Here is an explicit fixed path identifying the added form with
standard split stabilization:
\[
 q_t(z,p)=-t|z|^2+t|p|^2-(1-t)z\cdot p,\qquad 0\leq t\leq1.
 \tag{36}
\]
For each coordinate pair the gradient matrix \(H_t\) satisfies
\[
 H_t=\begin{pmatrix}-2t&t-1\\t-1&2t\end{pmatrix},\qquad
 H_t^2=(5t^2-2t+1)I
   =\left(5(t-\tfrac15)^2+\tfrac45\right)I.
 \tag{37}
\]
It has one positive and one negative eigenvalue, since its trace
is zero and its determinant is strictly negative. Thus \(q_t\)
has signature \((d,d)\) and a uniform gradient margin at all times.
Block orthogonality shows that \(g_o(X,w)+q_t(z,p)\) has no
nonzero critical point. The path is independent of \(X\).
Together with smoothing, it preserves the actual free stable tube
class on every fixed parameter slice, not just a numerical index
or a spherical image of that class.

The endpoint homogeneous function is smooth off the total origin
and over \(\overline U\). Compactness of the angular sphere and
homogeneity give new constants \(c_n>0,C_n<\infty\) as in (14),
and smooth extension provides uniform bounds on compact vertical
sets. This proves the analytic induction.
At the next step only (15), the smoothing parameter, and the
smoothing radius change. All are available for any finite
\(c_n,C_n\); they impose no smallness condition on the time increment.
The fixed finite number \(s\) from (13) therefore really suffices.

\section{The endpoint Hessian and composition of the accepted theorems}

After these \(s\) steps, the vertical dimension is
\[
 2m+2ds=2p,\qquad p=m+ds\geq20.
 \tag{38}
\]
The regular critical locus of \(F^{(s)}\) is diffeomorphically
parametrized onto \(\widetilde\varphi_1(\mathcal L_0)\).
By (7), over \(U_c\) its projection to the base is a diffeomorphism.
It follows that there is exactly one critical point over each
\(X\in U_c\), and that these points form a smooth section.

Let \(H=(F^{(s)})_{WW}\) at such a point. If \(Hv=0\), then
\[
 D(F^{(s)}_W)(0,v)=Hv=0.
 \tag{39}
\]
Full regularity identifies this kernel with the tangent space to
the critical locus. Therefore \((0,v)\) is in the kernel of the
differential of its base projection. That projection is a
diffeomorphism, so \(v=0\). Since \(H\) is square it is invertible.
Both hypotheses matter: \(F(X,w)=w^4\) has the zero graph and a
unique critical section, but has neither full regularity nor
nondegenerate vertical Hessians. This is not a counterexample to
the assertion with its actual hypotheses.

Restrict now to the fixed, unchanged embedding
\[
 S^9\longrightarrow U_c,\qquad b\longmapsto((b,v_0),0).
 \tag{40}
\]
The resulting \(f_b\) is smooth, has its one nondegenerate critical
section, and agrees with smooth homogeneous data outside one
uniform ball. By (36)--(38) the data have balanced type \(T_{p,p}\).
At each step their free stable class on (40) was preserved.
Initially it was \(a\), and the stable spherical class was zero;
the same is true at the endpoint.

A nonzero based sphere class cannot become freely null.
Indeed a path moving the basepoint identifies the two based classes
up to the fundamental-group action, which is an automorphism and
fixes zero. This applies also to the basepoint motion during the
regular homogeneous homotopies above. Consequently the endpoint
class is still freely non-null. The accepted one-critical theorem,
with exactly the hypotheses just verified, says it is freely null.
This contradiction proves Corollary~\ref{cor:main}.

\section{Audit and evidence boundary}

The reverse sign, Cartan action sign, mixed-generator signs, and
all four displayed cutoff identities of the manuscript's
\texttt{transport.tex} are consistent with (1).
The potential ambiguity in auxiliary momentum outside the spatial
support is resolved by (22), not by asserting \(p=p_o\) there.
Equations (24)--(25), (34)--(35), (13)--(15), and (36)--(37)
make the regularity, full-vertical error, finite-subdivision, and
fixed-stabilization arguments explicit. This is a verification of
the direct transport route, not a replacement invoking AIS.
No bypassed AIS lemma or unexamined Legendrian statement is certified.

The preceding accepted result established existence of \(L\),
but did not establish its non-isotopy. The present implication
is the missing link from that existence result to the already
accepted obstruction. It does not reconstruct either accepted proof.
The end-to-end conclusion retains their imported topological
assumptions and the separate pending formal citation-archive boundary.

The standalone \texttt{verify.py} and \texttt{exact\_certificate.json}
check the new signs and derivative identities, the full-regularity
block identity, the logarithmic-cutoff and error-budget algebra,
the fixed split path, and exact adversarial local examples.
The successful retained Jacobian call checks a symbolic scalar-jet instance of
the block product, with arbitrary old fiber Hessian, independently
of the SymPy replay. The general block argument is also stated
dimension-independently in (24)--(25).
Neither computation proves smooth ODE existence, the global
critical-locus correspondence, or the topology of the accepted
inputs. Those are supported by the mathematical proof, not
mislabelled as finite experiments or kernel formalization.
An earlier, equivalent input with denominator \(1+t\) was rejected
by the adapter's admission bound. That exact payload and structured
error are retained as a non-conclusion. The successful input names
the actual invertible block \(j=1+\sigma_{pX}\), as in (24).
There is no new external literature dependency and no novelty claim.
Independent Host-invoked review of this transport proof is still required.
\end{document}
